\chapter{Pleurodesis}

\section*{Overview}
Pleurodesis is a procedure that obliterates the pleural space by creating adhesions between the visceral and parietal pleura, preventing recurrent pneumothorax or pleural effusion.

\section*{Alternative Names}
Chemical pleurodesis; talc pleurodesis; pleural adhesion induction

\section*{Principle}
A sclerosant (most commonly sterile talc) is instilled into the pleural space, causing inflammation that fuses the two pleural layers together and eliminates the space in which air or fluid can accumulate.

\section*{History}
Pleurodesis was introduced in the early 20th century and refined with the use of talc poudrage and slurry in the mid-20th century.

\section*{Why It Is Done}
Performed for recurrent pneumothorax and for recurrent or malignant pleural effusion that recurs despite repeated drainage, in patients whose lungs can re-expand.

\section*{Is This a Primary or Secondary Test or Procedure}
Secondary treatment for recurrent pleural disease after initial drainage fails.

\section*{How It Is Performed}
Performed at the bedside or in the operating room. A chest tube drains the fluid or air, the lung is confirmed expanded, and a slurry of talc (or another agent) is instilled through the tube, which is clamped for a period while the patient changes position. In thoracoscopic (VATS) poudrage, talc is sprayed under vision.

\section*{Preparation}
Chest imaging to confirm lung expansion and absence of a trapped lung, analgesia, and informed consent.

\section*{Contraindications and Precautions}
An unexpanded (trapped) lung limits the surface contact and success of the procedure. Precaution: significant pain is expected and managed with analgesia.

\section*{Risks}
Fever and chest pain, infection, empyema, re-expansion pulmonary edema, and the rare risk of acute respiratory distress syndrome with talc.

\section*{Aftercare and Recovery}
The tube remains on suction until output is low, then is removed. Imaging confirms pleural apposition.

\section*{Outcome and Follow-up}
Successful pleurodesis is confirmed by lung expansion and lack of re-accumulation of air or fluid after tube removal.

\section*{Expected Results}
A fused pleural space with no recurrence of pneumothorax or effusion.

\section*{Follow-up}
Serial chest imaging for recurrence, and oncologic follow-up for malignant effusions.

\section*{When to Seek Medical Care}
Follow the procedure team's specific instructions. Seek urgent medical assessment for severe or worsening pain, uncontrolled bleeding, fever or signs of infection, trouble breathing, chest pain, fainting, new weakness or confusion, inability to urinate, or any rapidly worsening symptom. The appropriate warning signs depend on the procedure and the patient's medical condition.

\section*{References}
\begin{enumerate}
  \item MedlinePlus. Pleurodesis. U.S. National Library of Medicine; 2025.
  \item Current Medical Diagnosis and Treatment. McGraw-Hill; 2025.
\end{enumerate}

\clearpage
